\subsection{Rules}\label{ss:rules}

A rule rewrites the demanded node into its children, displayed on the full tuple; a component reappears
unchanged unless the rule modifies it. Conjunction is deterministic, disjunction branches (OR), and $\AX$
moves its argument into the context:
\[
(\wedge)\ \ \frac{\tup{\Gamma\cup\{\varphi\land\psi\}}{\Ctx}{Z}{\obligations}{\eventualities}}
{\tup{\Gamma\cup\{\varphi,\psi\}}{\Ctx}{Z}{\obligations}{\eventualities}}
\qquad
(\AX)\ \ \frac{\tup{\Gamma\cup\{\AX\varphi\}}{\Ctx}{Z}{\obligations}{\eventualities}}
{\tup{\Gamma}{\Ctx\cup\{\varphi\}}{Z}{\obligations}{\eventualities}}
\]
\[
(\vee)\ \ \frac{\tup{\Gamma\cup\{\varphi\lor\psi\}}{\Ctx}{Z}{\obligations}{\eventualities}}
{\tup{\Gamma\cup\{\varphi\}}{\Ctx}{Z}{\obligations}{\eventualities}\quad\big|\quad
\tup{\Gamma\cup\{\psi\}}{\Ctx}{Z}{\obligations}{\eventualities}}
\]
Here $\EX\psi$ and $\AX\psi$ are not modalities of the logic but {markers} that schedule an obligation at
the successors of a node. $\EX\psi$ asks that {one} successor satisfy $\psi$ and is discharged by
successor generation, which creates it; $\AX\psi$ asks that {every} successor satisfy $\psi$ and is
discharged by the context $\Ctx$, which rule $(\AX)$ fills and successor generation imposes on all children.
A pursued existential modality is propagated by $\EX$ and a universal one by $\AX$; we write $\nxQ\mu^\bullet$
for ``$\EX\mu^\bullet$ added to $\Gamma$ if $\mathrm{Q}=\mathrm{E}$, and $\AX\mu^\bullet$ added to $\Ctx$ if
$\mathrm{Q}=\mathrm{A}$''. For a pursuit, ``successor'' covers time-successors as well as discrete ones, so
the formula-clock advances by delay (made precise in successor generation below).

The two timed operators are discharged by one mechanism. The difficulty it resolves is that the deadline of a
modality $\mu=\mathrm{Q}(\varphi\,\odot_\kI\,\psi)$ ($\odot\in\{\mathbf{U},\mathbf{R}\}$) is measured {from
the configuration that demands it}; under a persistent context such as $\mathrm{AG}(\mu)$ the modality is
re-demanded at a {continuum} of configurations, each with its own absolute window. These demands cannot
be merged onto one running clock, nor afforded one clock each; instead each is verified {independently}
against a single formula-clock $x_\mu$, reset where the demand arises. The check does this in three movements:
it {launches} a pursuit of $\mu$ that runs $x_\mu$, {unfolds} the pursuit into a satisfaction set
$\Sat(\mu)$, and {decides} the demand against it. We describe the pursuit and its set, then the decision.

Write $\mu^\bullet$ for $\mu$ {pursued}. One unfolding step rewrites $\mu^\bullet$ as $x_\mu$ advances.
For until it either {stops} -- reaching $\psi$ inside the window, which fulfils the eventuality and so
discharges $\mu$ from $\eventualities$ -- or {goes} -- asserting $\varphi$ now and deferring the pursuit
to successors via $\nxQ\mu^\bullet$:
\[
\frac{\tup{\Gamma\cup\{\mu^\bullet\}}{\Ctx}{Z}{\obligations}{\eventualities}}
{\begin{array}{ll}
\tup{\Gamma\cup\{\psi\}}{\Ctx}{Z\cap(x_\mu\!\in\!\kI)}{\obligations}{\eventualities\setminus\{\mu\}} & (\text{stop})\\[4pt]
\tup{\Gamma\cup\{\varphi,\,\nxQ\mu^\bullet\}}{\Ctx}{Z}{\obligations}{\eventualities} & (\text{go})
\end{array}}
\qquad \mu=\mathrm{Q}(\varphi\,\mathbf{U}_\kI\,\psi)
\tag{$\mathbf{U}^\bullet$}\label{tag:Uunf}
\]
For release the pursuit partitions the zone by where $x_\mu$ lies relative to the window into its
before/in/after parts: $Z^{\in\kI}=Z\cap(x_\mu\in\kI)$, and $Z^{<\kI}$, $Z^{>\kI}$ the configurations below
resp.\ above the window, with an open boundary hyperplane (where $\kI$ excludes an endpoint) falling to the
adjacent part so the three genuinely partition $Z$ (the after-part is dropped when $\sup\kI=\infty$):
\[
\frac{\tup{\Gamma\cup\{\mu^\bullet\}}{\Ctx}{Z}{\obligations}{\eventualities}}
{\begin{array}{ll}
\tup{\Gamma\cup\{\varphi\}}{\Ctx}{Z^{<\kI}}{\obligations}{\eventualities} & (\text{before: }\varphi\text{ releases})\\[2pt]
\tup{\Gamma\cup\{\nxQ\mu^\bullet\}}{\Ctx}{Z^{<\kI}}{\obligations}{\eventualities} & (\text{before: continue})\\[2pt]
\tup{\Gamma\cup\{\psi,\,\nxQ\mu^\bullet\}}{\Ctx}{Z^{\in\kI}}{\obligations}{\eventualities} & (\text{in: }\psi,\text{ continue})\\[2pt]
\tup{\Gamma\cup\{\varphi,\psi\}}{\Ctx}{Z^{\in\kI}}{\obligations}{\eventualities} & (\text{in: }\psi,\ \varphi\text{ releases})\\[2pt]
\tup{\Gamma}{\Ctx}{Z^{>\kI}}{\obligations}{\eventualities} & (\text{after: lapsed})
\end{array}}
\qquad \mu=\mathrm{Q}(\varphi\,\mathbf{R}_\kI\,\psi)
\tag{$\mathbf{R}^\bullet$}\label{tag:Runf}
\]
Until needs only two cases -- a single in-window witness suffices -- whereas release, universal over the whole
window, needs the full before/in/after partition; the cases carrying $\nxQ\mu^\bullet$ are the recursive ones,
the rest base. An arbitrarily early $\varphi$ -- the before-window release case -- discharges the {entire}
window obligation, since it precedes every window position; this case is what lets a release hold with $\psi$
false throughout $\kI$, and it is necessary for completeness. A universal \emph{until} additionally caps its
go child (for $\mathrm{Q}=\mathrm{A}$) with the deadline invariant $x_\mu\le\sup\kI$ -- strict
($x_\mu<\sup\kI$) when $\kI$ excludes its supremum, so the cap matches the window and the stop edge stays
enabled up to it -- forcing the eventuality to a witness inside the window; a release carries no deadline -- it
must be free to lapse through the after-case -- and existential pursuits impose none. Repeated
unfolding generates finitely many pursuit nodes modulo $\cong_M$; propagating the stop / in-window base zones
backward along the go / continue edges -- to the least fixpoint for until, the greatest for release -- yields
the {satisfaction set}
\[
\Sat(\mu)\;=\;\{\nu : \text{$\mu$ holds at $\nu$ under the ambient context}\},
\qquad\text{a finite union of zones.}
\]
$\Sat(\mu)$ is an {operational} object: no model is presupposed, and its agreement with the semantics is
Lemma~\ref{lem:sat}.

A demanded $\mu$ on zone $Z$ is now discharged where $\Sat(\mu)$ covers it. The check {launches} the
pursuit just described -- resetting $x_\mu$, recording its clock in $\obligations$, and (for an until)
registering $\mu$ in $\eventualities$ -- and, when the cover is total, continues with $\mu$ met:
\[
\frac{\tup{\Gamma\cup\{\mu\}}{\Ctx}{Z}{\obligations}{\eventualities}}
{\begin{array}{ll}
\tup{\Gamma}{\Ctx}{Z}{\obligations}{\eventualities} & (\text{discharge})\\[3pt]
\tup{\Gamma\cup\{\mu^\bullet\}}{\Ctx}{Z\cap(x_\mu{=}0)}{\obligations\cup\{x_\mu\}}{\eventualities'} & (\text{launch})
\end{array}}
\qquad Z\subseteq\Sat(\mu)
\tag{\textrm{Chk}}\label{tag:chk}
\]
with $\eventualities'=\eventualities\cup\{\mu\}$ for an until and $\eventualities$ for a release. The two
children are conjunctive: the left continues the derivation with $\mu$ discharged, the right {is} the
pursuit, whose branches the witness model later reads its runs from (\S\ref{sec:correct}). When the cover is
only partial the demand is first refined to it, the rest closing,
\[
\frac{\tup{\Gamma\cup\{\mu\}}{\Ctx}{Z}{\obligations}{\eventualities}}
{\tup{\Gamma\cup\{\mu\}}{\Ctx}{Z_i}{\obligations}{\eventualities}}
\qquad Z_i\in\mathsf{cc}\big(Z\cap\Sat(\mu)\big)
\tag{\textrm{Chk}$^{*}$}\label{tag:chkstar}
\]
-- one disjunctive child per connected zone-component of $Z\cap\Sat(\mu)$ (each then dischargeable by
$(\mathrm{Chk})$), while the residual $Z\setminus\Sat(\mu)$ has no covering child and closes. The eventuality
of an until lives on the launch branch, where the loop condition of \S\ref{ss:loop} fulfils it; the discharged
continuation never carries it, so a spine loop answers only to the greatest-fixpoint (release) obligations it
retains -- a discharged until cannot stall a cycle. And because $(\mathrm{Chk}^{*})$ asserts $\mu$ only where
$\Sat(\mu)$ covers it, a universal demand higher up -- an $\mathrm{AG}$ needing {every} onset -- itself
closes when the cover is partial, so refinement is sound under any context.

Two features make this sound and complete. The retained context constrains witnesses by ambient safety, so an
$\mathrm{AG}_{\kI'}\neg q$ still forbids $q$ on them; and the universality $Z\subseteq\Sat(\mu)$ forces a
witness from {all} onsets of a persistent demand, not merely one -- a single onset that succeeds while a
neighbour fails would otherwise pass undetected. A single clock per modality therefore suffices, $K=O(n)$, the
continuum of onsets being absorbed symbolically into the zone $\Sat(\mu)$: this is the region quotient behind
the decidability of the fragment~\cite{fsttcs2000-LN,LaTorreNapoli2001TimedTreeAutomata}, internalised on the
fly. Untimed modalities are the case $\kI=[0,\infty)$, where the window constraint is vacuous and the check
reduces to ordinary reachability (until) or safety (release).

Finally, a saturated node is expanded by time-elapse and discrete successors. Time elapse advances $Z$ towards
$\uparrow Z$, {split at the integer hyperplanes} $x=b$ ($b\le M$) of active bounds: the procedure
advances only to the next such hyperplane, emitting an intermediate {time-successor} that inherits the
node's goals and context. Discrete successors realise the markers: for each $\EX\psi$ in $\Gamma$ one with
goals $\{\psi\}\cup\Ctx$, or, when none is demanded, a stutter successor with goals $\Ctx$. A pursuit's
residual $\nxQ\mu^\bullet$ is thus carried to {both} kinds of successor -- existentially to one
(\,$\mathrm{Q}=\mathrm{E}$, the marking's choice), universally to all (\,$\mathrm{Q}=\mathrm{A}$, via the
context) -- so $x_\mu$ advances by elapsing time as well as across edges. A timed window is therefore entered
by delay, the in-window stop is taken at the first successor that satisfies it, and a pure-delay witness --
waiting in one location until $x_\mu\in\kI$, then stopping -- is a chain of time-successors ending in a stop,
the stutter step supplying a discrete anchor only when the location must be left. Region stability
(\S\ref{sec:correct}) guarantees no in-window stop is skipped between two hyperplanes, and each piece being a
zone keeps the construction finite.





\paragraph{Timed Operators.}

Both timed operators are treated uniformly, by one formula-clock per modality and a single unfolding relation. Fix a modality $\mu=\mathrm{Q}(\varphi\,\odot_\kI\,\psi)$ with $\mathrm{Q}\in\{\mathrm{E},\mathrm{A}\}$ and $\odot\in\{\mathbf{U},\mathbf{R}\}$. The bound $\kI$ constrains the time elapsed from the configuration at which $\mu$ is demanded. For example, in case of a persistent context such as $\mathrm{AG}_{\kI '}(\mu)$, the demand recurs at a continuum of configurations over the $\kI'$, each with its own window \kI. 

To handle a demand $\mu$, we \emph{pursue} it. That is, we instantiate a formula-clock $x_\mu$ at the configuration that demands $\mu$, and use it to measure the elapsed time relevant to $\kI$. A single clock suffices for all onsets of a persistent demand, since these onsets are represented symbolically as a zone. We write $\mu^\bullet$ for the pursued modality, and $\nxQ\mu^\bullet$ for its propagation to successors: $\EX\mu^\bullet$ is added to $\Gamma$ when $\mathrm{Q}=\mathrm{E}$, and $\AX\mu^\bullet$ to the context $\Ctx$ when $\mathrm{Q}=\mathrm{A}$. An existential pursuit is thus carried as a goal and a universal pursuit through the context. In both cases $x_\mu$ advances at every successor, including time-successors.

%a formula-clock $x_\mu$, reset at the demanding configuration, measures the elapsed time relevant to $\kI$. A single clock suffices for all onsets of a persistent demand, since these onsets are represented symbolically as a zone rather than by one clock each. We write $\mu^\bullet$ for the pursued modality and $\nxQ\mu^\bullet$ for its propagation to successors: $\EX\mu^\bullet$ is added to $\Gamma$ when $\mathrm{Q}=\mathrm{E}$, and $\AX\mu^\bullet$ to the context $\Ctx$ when $\mathrm{Q}=\mathrm{A}$. An existential pursuit is thus carried as a goal and a universal pursuit through the context. In both cases $x_\mu$ advances at every successor, including time-successors.



For any until demand, i.e. of the form $\mu=\mathrm{Q}(\varphi\,\mathbf{U}_\kI\,\psi)$, the pursuit either {stops}, which happens when $\psi$ is satisfied inside the interval, or {goes}, which happens when $\varphi$ is satisfied and the pursuit continues, deferred via $\nxQ\mu^\bullet$.


%One unfolding step rewrites $\mu^\bullet$ as $x_\mu$ advances. For until it either {stops} -- reaching $\psi$ inside the window, which fulfils the eventuality and removes $\mu$ from $\eventualities$ -- or {goes}, asserting $\varphi$ now and deferring via $\nxQ\mu^\bullet$:
\[
\frac{\tup{\Gamma\cup\{\mu^\bullet\}}{\Ctx}{Z}{\obligations}{\eventualities}}
{\begin{array}{ll}
\tup{\Gamma\cup\{\psi\}}{\Ctx}{Z\cap(x_\mu\!\in\!\kI)}{\obligations}{\eventualities\setminus\{\mu\}} & (\text{stop})\\[4pt]
\tup{\Gamma\cup\{\varphi,\,\EX\mu^\bullet\}}{\Ctx}{Z}{\obligations}{\eventualities} & (\text{go},\ \mathrm{Q}{=}\mathrm{E})\\[4pt]
\tup{\Gamma\cup\{\varphi\}}{\Ctx\cup\{\AX\mu^\bullet\}}{Z\cap(x_\mu\le\sup\kI)}{\obligations}{\eventualities} & (\text{go-A},\ \mathrm{Q}{=}\mathrm{A})
\end{array}}
\qquad \mu=\mathrm{Q}(\varphi\,\mathbf{U}_\kI\,\psi)
\tag{$\mathbf{U}^\bullet$}\label{tag:Uunf}
\]

For release demands, i.e. of the form $\mu=\mathrm{Q}(\varphi\,\mathbf{R}_\kI\,\psi)$, we partition the zone by where $x_\mu$ lies relative to \kI. This results in three cases, which we call the {before}, {in}, and {after} parts of the \kI. Write $Z^{\in\kI}=Z\cap(x_\mu\in\kI)$ for the configurations inside \kI, and $Z^{<\kI}$, $Z^{>\kI}$ for those below and above \kI, respectively. An open boundary hyperplane (where $\kI$ excludes an endpoint) falls to the adjacent part so that the three genuinely partition $Z$; the after-part is dropped when $\sup\kI=\infty$.


%

%into its before/in/after parts: $Z^{\in\kI}=Z\cap(x_\mu\in\kI)$, and $Z^{<\kI}$, $Z^{>\kI}$ the configurations below resp.\ above the window, with an open boundary hyperplane (where $\kI$ excludes an endpoint) falling to the adjacent part so the three genuinely partition $Z$ (the after-part is dropped when $\sup\kI=\infty$):


\[
\frac{\tup{\Gamma\cup\{\mu^\bullet\}}{\Ctx}{Z}{\obligations}{\eventualities}}
{\begin{array}{ll}
\tup{\Gamma\cup\{\varphi\}}{\Ctx}{Z^{<\kI}}{\obligations}{\eventualities} & (\text{before: }\varphi\text{ releases})\\[2pt]
\tup{\Gamma\cup\{\nxQ\mu^\bullet\}}{\Ctx}{Z^{<\kI}}{\obligations}{\eventualities} & (\text{before: continue})\\[2pt]
\tup{\Gamma\cup\{\psi,\,\nxQ\mu^\bullet\}}{\Ctx}{Z^{\in\kI}}{\obligations}{\eventualities} & (\text{in: }\psi,\text{ continue})\\[2pt]
\tup{\Gamma\cup\{\varphi,\psi\}}{\Ctx}{Z^{\in\kI}}{\obligations}{\eventualities} & (\text{in: }\psi,\ \varphi\text{ releases})\\[2pt]
\tup{\Gamma}{\Ctx}{Z^{>\kI}}{\obligations}{\eventualities} & (\text{after: lapsed})
\end{array}}
\qquad \mu=\mathrm{Q}(\varphi\,\mathbf{R}_\kI\,\psi)
\tag{$\mathbf{R}^\bullet$}\label{tag:Runf}
\]



The two rules differ in shape for a structural reason. An until asserts an eventuality, and a single position inside \kI witnesses it; the stop case records that witness, so three cases suffice. A release is universal over the \kI, and no single position discharges it, so the pursuit splits the \kI into its before, in, and after parts, giving five cases. 

%In both rules the cases carrying $\nxQ\mu^\bullet$ are recursive, that is they defer the pursuit to a successor, and the remaining cases are base, that is they discharge, or settle the pursuit. 

%Two properties of the base cases are worth stating. First, in a release an occurrence of $\varphi$ strictly before the window discharges the entire window obligation, because it precedes every position of $\kI$; this is why a release may hold with $\psi$ false throughout $\kI$. Second, a universal until ($\mathrm{Q}=\mathrm{A}$) defers through the go-A case, which adds the deadline invariant $x_\mu\le\sup\kI$ (strict when $\kI$ excludes its supremum) to the zone. Time then cannot elapse past the window while the eventuality is pending, forcing a witness inside $\kI$. An existential until defers through the go case, which imposes no bound; a release imposes none either, since it must remain free to reach its after-case and lapse.


Unfolding $\mu^\bullet$ produces two kinds of child node: (i) \emph{base children}, which settle the demand at the current configuration and do not carry $\mu^\bullet$ onward (the stop case of an until, and the three discharging cases of a release); and (ii) \emph{recursive children}, which instead defer $\nxQ\mu^\bullet$ to their successors, where $\mu$ must be pursued again (the go and go-A cases of an until, and the two continue cases of a release). Repeated unfolding thus determines a set of clock valuations at which $\mu$ holds, which we call $\Sat(\mu)$. Let $B_\mu$ be the union of the base-case zones. A valuation $\nu$ lies in $\Sat(\mu)$ exactly when it lies in $B_\mu$, or it matches a recursive child whose deferred successors lie in $\Sat(\mu)$ (some successor when $\mathrm{Q}=\mathrm{E}$, all of them when $\mathrm{Q}=\mathrm{A}$). This is a recursive equation for $\Sat(\mu)$. Enumerating valuations directly would not terminate (as there are uncountably many), but because zones are compared up to $\cong_M$ the unfolding reaches only finitely many distinct nodes, so the equation is solved by iterating over a finite family of zones. It has a least and a greatest solution, and the two cases take opposite ones. For an until, $\Sat(\mu)$ is the least: a valuation qualifies only if it can \emph{reach} a base case in finitely many steps. For a release, it is the greatest: a valuation qualifies as long as it can continue without ever violating the release. In either case $\Sat(\mu)$ is a finite union of zones, assembled from the zones of the base and recursive children, and computed from the rules alone with no model assumed.

%In either case $\Sat(\mu)$ is computed using union of finitely many zones, which terminates.  

%In either case $\Sat(\mu)$ is a finite union of zones, computed from the rules alone with no model assumed.







\paragraph{Deciding a demand: the cover rule.}
A demanded $\mu$ on zone $Z$ holds exactly where $\Sat(\mu)$ covers it. Rule $(\mathrm{Cov})$ {launches} the pursuit above -- resetting $x_\mu$, recording its clock in $\obligations$, and (for an until) registering $\mu$ in $\eventualities$ -- and, when the cover is total, continues with $\mu$ met:
\[
\frac{\tup{\Gamma\cup\{\mu\}}{\Ctx}{Z}{\obligations}{\eventualities}}
{\begin{array}{ll}
\tup{\Gamma}{\Ctx}{Z}{\obligations}{\eventualities} & (\text{discharge})\\[3pt]
\tup{\Gamma\cup\{\mu^\bullet\}}{\Ctx}{Z\cap(x_\mu{=}0)}{\obligations\cup\{x_\mu\}}{\eventualities\cup\{\mu\}} & (\text{launch-U})\\[3pt]
\tup{\Gamma\cup\{\mu^\bullet\}}{\Ctx}{Z\cap(x_\mu{=}0)}{\obligations\cup\{x_\mu\}}{\eventualities} & (\text{launch-R})
\end{array}}
\ Z\subseteq\Sat(\mu)
\tag{\textrm{Cov}}\label{tag:chk}
\]
The two children are conjunctive: the left continues the derivation with $\mu$ discharged, the right {is} the pursuit, kept as the witness whose branches supply the model's runs. When the cover is only partial the demand is first refined to it, the rest closing:
\[
\frac{\tup{\Gamma\cup\{\mu\}}{\Ctx}{Z}{\obligations}{\eventualities}}
{\tup{\Gamma\cup\{\mu\}}{\Ctx}{Z_i}{\obligations}{\eventualities}}
\qquad Z_i\in\mathsf{cc}\big(Z\cap\Sat(\mu)\big)
\tag{\textrm{Cov}$^{*}$}\label{tag:chkstar}
\]
one disjunctive child per connected zone-component of $Z\cap\Sat(\mu)$ (each then dischargeable by $(\mathrm{Cov})$), while the residual $Z\setminus\Sat(\mu)$ has no covering child and closes. The eventuality of an until rides only the launch branch; the discharged continuation never carries it, so a discharged until cannot stall a cycle. Untimed modalities are the case $\kI=[0,\infty)$: the window constraint is vacuous and the rule reduces to plain reachability (until) or safety (release).






Both timed operators are treated uniformly, by one formula-clock per modality and a single unfolding relation. Fix a modality $\mu=\mathrm{Q}(\varphi\,\odot_\kI\,\psi)$ with $\mathrm{Q}\in\{\mathrm{E},\mathrm{A}\}$ and $\odot\in\{\mathbf{U},\mathbf{R}\}$. The bound $\kI$ constrains the time elapsed from the configuration at which $\mu$ is demanded. For example, in case of a persistent context such as $\mathrm{AG}_{\kI '}(\mu)$, the demand recurs at a continuum of configurations over the $\kI'$, each with its own window \kI. 

To handle a demand $\mu$, we \emph{pursue} it. That is, we instantiate a formula-clock $x_\mu$ at the configuration that demands $\mu$, and use it to measure the elapsed time relevant to $\kI$. A single clock suffices for all onsets of a persistent demand, since these onsets are represented symbolically as a zone. We write $\mu^\bullet$ for the pursued modality, and $\nxQ\mu^\bullet$ for its propagation to successors: $\EX\mu^\bullet$ is added to $\Gamma$ when $\mathrm{Q}=\mathrm{E}$, and $\AX\mu^\bullet$ to the context $\Ctx$ when $\mathrm{Q}=\mathrm{A}$. An existential pursuit is thus carried as a goal and a universal pursuit through the context. In both cases $x_\mu$ advances at every successor, including time-successors.

%a formula-clock $x_\mu$, reset at the demanding configuration, measures the elapsed time relevant to $\kI$. A single clock suffices for all onsets of a persistent demand, since these onsets are represented symbolically as a zone rather than by one clock each. We write $\mu^\bullet$ for the pursued modality and $\nxQ\mu^\bullet$ for its propagation to successors: $\EX\mu^\bullet$ is added to $\Gamma$ when $\mathrm{Q}=\mathrm{E}$, and $\AX\mu^\bullet$ to the context $\Ctx$ when $\mathrm{Q}=\mathrm{A}$. An existential pursuit is thus carried as a goal and a universal pursuit through the context. In both cases $x_\mu$ advances at every successor, including time-successors.



For any until demand, i.e. of the form $\mu=\mathrm{Q}(\varphi\,\mathbf{U}_\kI\,\psi)$, the pursuit either {stops}, which happens when $\psi$ is satisfied inside the interval, or {goes}, which happens when $\varphi$ is satisfied and the pursuit continues, deferred via $\nxQ\mu^\bullet$.


%One unfolding step rewrites $\mu^\bullet$ as $x_\mu$ advances. For until it either {stops} -- reaching $\psi$ inside the window, which fulfils the eventuality and removes $\mu$ from $\eventualities$ -- or {goes}, asserting $\varphi$ now and deferring via $\nxQ\mu^\bullet$:
\[
\frac{\tup{\Gamma\cup\{\mu^\bullet\}}{\Ctx}{Z}{\obligations}{\eventualities}}
{\begin{array}{ll}
\tup{\Gamma\cup\{\psi\}}{\Ctx}{Z\cap(x_\mu\!\in\!\kI)}{\obligations}{\eventualities\setminus\{\mu\}} & (\text{stop})\\[4pt]
\tup{\Gamma\cup\{\varphi,\,\EX\mu^\bullet\}}{\Ctx}{Z}{\obligations}{\eventualities} & (\text{go},\ \mathrm{Q}{=}\mathrm{E})\\[4pt]
\tup{\Gamma\cup\{\varphi\}}{\Ctx\cup\{\AX\mu^\bullet\}}{Z\cap(x_\mu\le\sup\kI)}{\obligations}{\eventualities} & (\text{go-A},\ \mathrm{Q}{=}\mathrm{A})
\end{array}}
\qquad \mu=\mathrm{Q}(\varphi\,\mathbf{U}_\kI\,\psi)
\tag{$\mathbf{U}^\bullet$}\label{tag:Uunf}
\]

For release demands, i.e. of the form $\mu=\mathrm{Q}(\varphi\,\mathbf{R}_\kI\,\psi)$, we partition the zone by where $x_\mu$ lies relative to \kI. This results in three cases, which we call the {before}, {in}, and {after} parts of the \kI. Write $Z^{\in\kI}=Z\cap(x_\mu\in\kI)$ for the configurations inside \kI, and $Z^{<\kI}$, $Z^{>\kI}$ for those below and above \kI, respectively. An open boundary hyperplane (where $\kI$ excludes an endpoint) falls to the adjacent part so that the three genuinely partition $Z$; the after-part is dropped when $\sup\kI=\infty$.


%

%into its before/in/after parts: $Z^{\in\kI}=Z\cap(x_\mu\in\kI)$, and $Z^{<\kI}$, $Z^{>\kI}$ the configurations below resp.\ above the window, with an open boundary hyperplane (where $\kI$ excludes an endpoint) falling to the adjacent part so the three genuinely partition $Z$ (the after-part is dropped when $\sup\kI=\infty$):


\[
\frac{\tup{\Gamma\cup\{\mu^\bullet\}}{\Ctx}{Z}{\obligations}{\eventualities}}
{\begin{array}{ll}
\tup{\Gamma\cup\{\varphi\}}{\Ctx}{Z^{<\kI}}{\obligations}{\eventualities} & (\text{before: }\varphi\text{ releases})\\[2pt]
\tup{\Gamma\cup\{\nxQ\mu^\bullet\}}{\Ctx}{Z^{<\kI}}{\obligations}{\eventualities} & (\text{before: continue})\\[2pt]
\tup{\Gamma\cup\{\psi,\,\nxQ\mu^\bullet\}}{\Ctx}{Z^{\in\kI}}{\obligations}{\eventualities} & (\text{in: }\psi,\text{ continue})\\[2pt]
\tup{\Gamma\cup\{\varphi,\psi\}}{\Ctx}{Z^{\in\kI}}{\obligations}{\eventualities} & (\text{in: }\psi,\ \varphi\text{ releases})\\[2pt]
\tup{\Gamma}{\Ctx}{Z^{>\kI}}{\obligations}{\eventualities} & (\text{after: lapsed})
\end{array}}
\qquad \mu=\mathrm{Q}(\varphi\,\mathbf{R}_\kI\,\psi)
\tag{$\mathbf{R}^\bullet$}\label{tag:Runf}
\]



The two rules differ in shape for a structural reason. An until asserts an eventuality, and a single position inside \kI witnesses it; the stop case records that witness, so three cases suffice. A release is universal over the \kI, and no single position discharges it, so the pursuit splits the \kI into its before, in, and after parts, giving five cases. 

%In both rules the cases carrying $\nxQ\mu^\bullet$ are recursive, that is they defer the pursuit to a successor, and the remaining cases are base, that is they discharge, or settle the pursuit. 

%Two properties of the base cases are worth stating. First, in a release an occurrence of $\varphi$ strictly before the window discharges the entire window obligation, because it precedes every position of $\kI$; this is why a release may hold with $\psi$ false throughout $\kI$. Second, a universal until ($\mathrm{Q}=\mathrm{A}$) defers through the go-A case, which adds the deadline invariant $x_\mu\le\sup\kI$ (strict when $\kI$ excludes its supremum) to the zone. Time then cannot elapse past the window while the eventuality is pending, forcing a witness inside $\kI$. An existential until defers through the go case, which imposes no bound; a release imposes none either, since it must remain free to reach its after-case and lapse.


Unfolding $\mu^\bullet$ produces two kinds of child node: (i) \emph{base children}, which settle the demand at the current configuration and do not carry $\mu^\bullet$ onward (the stop case of an until, and the three discharging cases of a release); and (ii) \emph{recursive children}, which instead defer $\nxQ\mu^\bullet$ to their successors, where $\mu$ must be pursued again (the go and go-A cases of an until, and the two continue cases of a release). 



Repeated unfolding thus determines a set of clock valuations at which $\mu$ holds, which we call $\Sat(\mu)$. Let $B_\mu$ be the union of the base-case zones. A valuation $\nu$ lies in $\Sat(\mu)$ exactly when it lies in $B_\mu$, or it matches a recursive child whose deferred successors lie in $\Sat(\mu)$ (some successor when $\mathrm{Q}=\mathrm{E}$, all of them when $\mathrm{Q}=\mathrm{A}$). This is a recursive equation for $\Sat(\mu)$. Enumerating valuations directly would not terminate (as there are uncountably many), but because zones are compared up to $\cong_M$ the unfolding reaches only finitely many distinct nodes, so the equation is solved by iterating over a finite family of zones. It has a least and a greatest solution, and the two cases take opposite ones. For an until, $\Sat(\mu)$ is the least: a valuation qualifies only if it can \emph{reach} a base case in finitely many steps. For a release, it is the greatest: a valuation qualifies as long as it can continue without ever violating the release. In either case $\Sat(\mu)$ is a finite union of zones, assembled from the zones of the base and recursive children, and computed from the rules alone with no model assumed.

%In either case $\Sat(\mu)$ is computed using union of finitely many zones, which terminates.  

%In either case $\Sat(\mu)$ is a finite union of zones, computed from the rules alone with no model assumed.







\paragraph{Deciding a demand: the cover rule.}
A demanded $\mu$ on zone $Z$ holds exactly where $\Sat(\mu)$ covers it. Rule $(\mathrm{Cov})$ {launches} the pursuit above -- resetting $x_\mu$, recording its clock in $\obligations$, and (for an until) registering $\mu$ in $\eventualities$ -- and, when the cover is total, continues with $\mu$ met:
\[
\frac{\tup{\Gamma\cup\{\mu\}}{\Ctx}{Z}{\obligations}{\eventualities}}
{\begin{array}{ll}
\tup{\Gamma}{\Ctx}{Z}{\obligations}{\eventualities} & (\text{discharge})\\[3pt]
\tup{\Gamma\cup\{\mu^\bullet\}}{\Ctx}{Z\cap(x_\mu{=}0)}{\obligations\cup\{x_\mu\}}{\eventualities\cup\{\mu\}} & (\text{launch-U})\\[3pt]
\tup{\Gamma\cup\{\mu^\bullet\}}{\Ctx}{Z\cap(x_\mu{=}0)}{\obligations\cup\{x_\mu\}}{\eventualities} & (\text{launch-R})
\end{array}}
\ Z\subseteq\Sat(\mu)
\tag{\textrm{Cov}}\label{tag:chk}
\]
The two children are conjunctive: the left continues the derivation with $\mu$ discharged, the right {is} the pursuit, kept as the witness whose branches supply the model's runs. When the cover is only partial the demand is first refined to it, the rest closing:
\[
\frac{\tup{\Gamma\cup\{\mu\}}{\Ctx}{Z}{\obligations}{\eventualities}}
{\tup{\Gamma\cup\{\mu\}}{\Ctx}{Z_i}{\obligations}{\eventualities}}
\qquad Z_i\in\mathsf{cc}\big(Z\cap\Sat(\mu)\big)
\tag{\textrm{Cov}$^{*}$}\label{tag:chkstar}
\]
one disjunctive child per connected zone-component of $Z\cap\Sat(\mu)$ (each then dischargeable by $(\mathrm{Cov})$), while the residual $Z\setminus\Sat(\mu)$ has no covering child and closes. The eventuality of an until rides only the launch branch; the discharged continuation never carries it, so a discharged until cannot stall a cycle. Untimed modalities are the case $\kI=[0,\infty)$: the window constraint is vacuous and the rule reduces to plain reachability (until) or safety (release).